\chapter{複合ボソン、スピン表現、ディラックの倍化}

\section{交換対称性と複合ボソン}

冒頭の問いは、複合粒子に関する真のパラドックスである。2つのフェルミオンがボソンを作るなら、そのようなボソン2つがフェルミオンの排他律に違反することなく同じ状態を占めることができるのか。答えは交換規則そのものから始まる。

2つの同種粒子が2つの引数をもつ波動関数で記述されるとしよう。記号\(x\)と\(y\)は2つの粒子の枠をラベルするものであり、2つの種を区別するものではない。ボソンでは波動関数は対称であり、フェルミオンでは反対称である：
\begin{equation}
\Psi_B(x,y)=\Psi_B(y,x),
\qquad
\Psi_F(x,y)=-\Psi_F(y,x).
\end{equation}
この対称性は、2つの識別不能な粒子を入れ替えても新たな物理状態は生じないことを述べている。

2つのボソンは確かに同じ一粒子状態を占めることができる。その状態の波動関数を\(\varphi\)とすると、2粒子の積は
\begin{equation}
\Psi(x,y)=\varphi(x)\varphi(y).
\end{equation}
波動関数は演算子や行列ではなく通常の複素数値関数であるから、その値は可換である：
\begin{equation}
\varphi(x)\varphi(y)=\varphi(y)\varphi(x).
\end{equation}
したがってこの積は対称である。同じ構成では非零の2フェルミオン状態は得られない。反対称性は
\begin{equation}
\varphi(x)\varphi(y)=-\varphi(y)\varphi(x),
\end{equation}
を要求するが、負号のない両辺は等しいからである。

より一般には、交換対称性を仮定しない任意の2粒子波動関数\(F(x,y)\)から始める。その対称部分と反対称部分は、規格化を除いて
\begin{align}
F_{\mathrm{sym}}(x,y)&=F(x,y)+F(y,x),\\
F_{\mathrm{asym}}(x,y)&=F(x,y)-F(y,x).
\end{align}
である。交換の検査は明示的である：
\begin{align}
F_{\mathrm{sym}}(y,x)
&=F(y,x)+F(x,y)
=F_{\mathrm{sym}}(x,y),\\
F_{\mathrm{asym}}(y,x)
&=F(y,x)-F(x,y)
=-F_{\mathrm{asym}}(x,y).
\end{align}
交換した関数を加えると対称化され、交換した関数を引くと反対称化される。

水素原子の対は、より興味深い場合を与える。4つの構成フェルミオン、すなわち2つの同種電子と2つの同種陽子がある。\(x_1,x_2\)で電子の変数を表し、\(y_1,y_2\)で陽子の変数を表すことにしよう。4体波動関数は、電子の交換のもとで、また陽子の交換のもとで、それぞれ別個に反対称でなければならない：
\begin{align}
\Psi(x_2,x_1;y_1,y_2)
&=-\Psi(x_1,x_2;y_1,y_2),\\
\Psi(x_1,x_2;y_2,y_1)
&=-\Psi(x_1,x_2;y_1,y_2).
\end{align}

2つの異なる一原子波動関数を動機づけるために、一方の水素原子が原点付近に局在し、波動関数\(\psi(x,y)\)をもつとし、他方が別の場所に局在し、波動関数\(\phi(x,y)\)をもつと想像しよう。後者は単に前者の平行移動されたコピーであってもよく、その電子と陽子の引数の両方を平行移動して得られるが、それでもこの設定では依然として異なる波動関数である。素朴な積は
\begin{equation}
\psi(x_1,y_1)\phi(x_2,y_2).
\end{equation}
であり、これはまだ要求されるフェルミオン的反対称性のいずれももっていない。

まず電子の変数について反対称化する：
\begin{equation}
\Psi_x=
\psi(x_1,y_1)\phi(x_2,y_2)
-\psi(x_2,y_1)\phi(x_1,y_2).
\end{equation}
\(x_1\leftrightarrow x_2\)のもとで、2つの項は入れ替わり、式全体は符号を変える。次にこの結果を陽子の変数について反対称化する。4項の順序は2回の逐次的反対称化によって固定される。\footnote{この構成はまた、口頭の計算中に行われた変数ラベルの修正を通じて符号パターン\(+,-,-,+\)を固定する。}
\begin{align}
\Psi_{2H}(x_1,x_2;y_1,y_2)
&=\Psi_x(x_1,x_2;y_1,y_2)\nonumber\\
&\quad{}-\Psi_x(x_1,x_2;y_2,y_1)\\
&=\psi(x_1,y_1)\phi(x_2,y_2)\nonumber\\
&\quad{}-\psi(x_2,y_1)\phi(x_1,y_2)\nonumber\\
&\quad{}-\psi(x_1,y_2)\phi(x_2,y_1)\nonumber\\
&\quad{}+\psi(x_2,y_2)\phi(x_1,y_1).
\label{eq:two-hydrogen-antisymmetrized}
\end{align}
電子の交換は最初の2項と最後の2項を入れ替え、全体の符号を反転させる。陽子の交換も同じことをする。したがってこの状態は、構成フェルミオンに要求されるちょうど2つの別個の反対称性をもつ。

\begin{classroomqa}
\lecturetimestamp{00:11:11}
\audiencequestion{水素原子がボソンなら、2つを直接重ねて置くことを妨げるものは何か。}
\lecturerresponse{ボソン的交換統計は原子間の力を取り除かない。構成要素間のパウリ排他律は、他の原子間力とともに、空間的重なりを極めて高コストにしうる。これは、原子が同じ非局在化した運動量状態を占めてよいかどうかを問うこととは異なる。運動量固有状態は空間全体に広がっているので、共通の運動量は原子を同じ点に衝突させることを意味しない。ビリヤード球のような硬い芯をもつ物体でさえボソン的交換統計に従いうる：その硬い芯は空間的重なりを妨げるが、ボソン統計は依然として共通の運動量占有を許すであろう。}
\end{classroomqa}

\begin{classroomqa}
\lecturetimestamp{00:12:28}
\audiencequestion{2つの水素原子を同じ運動量固有状態に置くと、その構成フェルミオンは同じ状態に置かれることになるのか。}
\lecturerresponse{いいえ。4体状態は、2つの電子変数の交換のもとで反対称のままであり、また2つの陽子変数の交換のもとで別個に反対称のままである。両方の構成要素の同時交換---2つの原子全体の交換---だけが対称である。}
\end{classroomqa}

実際、二つの原子を交換することは、両方の置換を行うことを意味する。
\begin{equation}
x_1\leftrightarrow x_2,
\qquad
y_1\leftrightarrow y_2.
\end{equation}
第一の置換はフェルミオン的なマイナス符号を一つ寄与し、第二の置換はさらに一つ寄与する：
\begin{equation}
\begin{aligned}
\Psi_{2H}(x_2,x_1;y_2,y_1)
&=(-1)(-1)\Psi_{2H}(x_1,x_2;y_1,y_2)\\
&=\Psi_{2H}(x_1,x_2;y_1,y_2).
\end{aligned}
\end{equation}
同じ結果は式~\eqref{eq:two-hydrogen-antisymmetrized}で項ごとに見ることもできる。原子全体の交換の下で、
\begin{equation}
\psi(x_1,y_1)\phi(x_2,y_2)
\longleftrightarrow
\psi(x_2,y_2)\phi(x_1,y_1),
\end{equation}
であるから、第一項と第四項が交換される。同様に、
\begin{equation}
-\psi(x_2,y_1)\phi(x_1,y_2)
\longleftrightarrow
-\psi(x_1,y_2)\phi(x_2,y_1),
\end{equation}
であるから、第二項と第三項が交換される。完全な状態は原子の交換の下で対称である。二つの構成要素のマイナス符号が合成系にプラス符号を与える。

この交換の言明は、原子を重ね合わせるのに必要なエネルギーとは論理的に独立である。強い反発によってそのような重なりが実現不可能になることはあるが、対称な交換法則は変わらない。

\begin{classroomqa}
\lecturetimestamp{00:14:42}
\audiencequestion{構成された四体状態は物理的に何を表すのか、そしてそれは特に運動量空間の状態なのか。}
\lecturerresponse{それは、二つの一体波動関数から、まず電子のラベルを反対称化し、次に陽子のラベルを反対称化して得られる二原子波動関数である。必ずしも運動量空間で書かれる必要はない。変数は位置、運動量、あるいは別の表現を記述してよい。要求される構成要素の交換対称性は同じである。}
\end{classroomqa}

したがって、二つの原子は同じ非局在化した運動量波動関数を占めることがあり得る。ただし、それらの空間分布を強い重なりに押し込むには多大なエネルギーを要することがある。

ここで二つの一体波動関数を等しいと置く：\(\phi=\psi\)。式~\eqref{eq:two-hydrogen-antisymmetrized}は次になる。
\begin{align}
\Psi_{\mathrm{same}}
&=\psi(x_1,y_1)\psi(x_2,y_2)
-\psi(x_2,y_1)\psi(x_1,y_2)\nonumber\\
&\quad
-\psi(x_1,y_2)\psi(x_2,y_1)
+\psi(x_2,y_2)\psi(x_1,y_1).
\end{align}

この四体状態は消えるか。消えない。第一項と第四項の積は等しく、第二項と第三項の積も等しい。一体波動関数の可換性を用いると、
\begin{equation}
\Psi_{\mathrm{same}}
=
2\left[
\psi(x_1,y_1)\psi(x_2,y_2)
-\psi(x_2,y_1)\psi(x_1,y_2)
\right],
\label{eq:same-hydrogen-state}
\end{equation}
となり、これは一般にゼロでない。それは二つの電子変数について反対称であり、二つの陽子変数について別個に反対称である。

ここには、それぞれ鋭く局在した原子に一つの鋭く局在した電子を割り当てることはない。原子に束縛された電子は、その原子波動関数にわたって空間的に非局在化している。したがって、二つの原子に同じ原子状態を与えても、\(x_1=x_2\)が成り立つわけでも、\(y_1=y_2\)が成り立つわけでもない。

構成要素の排他性は、一致を直接調べることで見える。モデル化された二つの電子引数が一致する場合、\(x_1=x_2=x\)、式~\eqref{eq:same-hydrogen-state}は次を与える。
\begin{align}
\Psi_{\mathrm{same}}(x,x;y_1,y_2)
&=2\bigl[
\psi(x,y_1)\psi(x,y_2)\nonumber\\
&\quad{}-\psi(x,y_1)\psi(x,y_2)
\bigr]\\
&=0.
\end{align}
対応する表式は\(y_1=y_2\)のときにも消える。したがって、同一状態の合成波動関数はゼロでないが、一致する同一フェルミオン引数におけるその振幅は消える。\footnote{このモデルではスピンおよびその他の一粒子ラベルは抑制されている。フェルミオン的反対称性は一粒子ラベルの完全な集合に適用される。表示された相殺は、ここで保持された引数の一致に関するものである。}

\begin{classroomqa}
\lecturetimestamp{00:20:41}
\audiencequestion{原子が同じ状態にありながら、それらの電子と陽子が同じ構成要素状態にないことはあり得るか。}
\lecturerresponse{はい。合成原子波動関数が等しいことは、電子と陽子の別個の反対称性を取り除かない。完全な四体状態は、二つの原子が一つの原子状態にあることを記述しながら、一致する同一フェルミオン引数にゼロ振幅を割り当てることができる。}
\end{classroomqa}

\begin{classroomqa}
\lecturetimestamp{00:21:12}
\audiencequestion{反対称化は、同一構成要素が一致する部分を射影によって取り除いたのか。}
\lecturerresponse{はい。ここに表示された変数では、それらの一致成分は正確に相殺する。正確な言明は、二つの同一フェルミオン引数が等しいとき反対称波動関数が消えるということである。}
\end{classroomqa}

この相殺は、まず原子波動関数を等しくすることに依存していなかった。式~\eqref{eq:two-hydrogen-antisymmetrized}の一般状態について、
\begin{align}
\Psi_{2H}(x,x;y_1,y_2)
&=\psi(x,y_1)\phi(x,y_2)\nonumber\\
&\quad{}-\psi(x,y_1)\phi(x,y_2)\nonumber\\
&\quad{}-\psi(x,y_2)\phi(x,y_1)\nonumber\\
&\quad{}+\psi(x,y_2)\phi(x,y_1)\\
&=0.
\end{align}

\begin{classroomqa}
\lecturetimestamp{00:21:32}
\audiencequestion{原子が異なる状態にあり、同一構成要素変数が一致するときにも、この相殺は成り立つか。}
\lecturerresponse{はい。それは\(\phi=\psi\)と置くことからではなく、反対称性そのものから従う。同じ言明は別の表現でも成り立つ。運動量空間に変換した後、二つの同一フェルミオン運動量引数が一致し、かつ他の抑制された一粒子ラベルも同様に一致するならば、振幅は消える。}
\end{classroomqa}

\section{有限次元量子系としてのスピン}

スピンの数学は、アイソスピン、カラー、および素粒子物理学のその他の内部対称性において繰り返し現れる。この類似は形式的である。同じ代数が現れるが、物理量が同じ解釈を持つ必要はない。当面の課題はスピンをもう少し展開し、その後ディラック方程式に戻ることである。この道具立ては後にアイソスピンに用いられる。アイソスピンは歴史的に1930年代に関連づけられた概念である。

電子はスピン\(\tfrac12\)の粒子であり、陽子やミューオンも同様である。そのスピンの\(z\)成分は2つの値をとりうる。単位を定める前に、
\begin{equation}
S_z=m\hbar.
\end{equation}
と書く。以後\(\hbar=1\)とする。文字\(L\)は慣例的に軌道角運動量に留保されているので、固有スピンは\(S\)で表す。これは同じ角運動量代数に従う：
\begin{equation}
[S_x,S_y]=iS_z,
\quad
[S_y,S_z]=iS_x,
\quad
[S_z,S_x]=iS_y,
\end{equation}
および
\begin{equation}
S_\pm=S_x\pm iS_y.
\end{equation}
演算子\(S_+\)は\(S_z\)固有値を上げ、\(S_-\)は下げる。

\begin{classroomqa}
\lecturetimestamp{00:25:39}
\audiencequestion{議論は特に固有スピンについてですか。}
\lecturerresponse{固有スピンが意図された焦点である。同じ角運動量代数はより一般に、軌道角運動量や巨視的物体の回転にも当てはまる。}
\end{classroomqa}

スピン1/2では、\(S_z\)の固有値は\(+\tfrac12\)と\(-\tfrac12\)だけである。今は電子についての他のすべてを抑え、その位置と運動量を忘れて2つのスピン状態だけを保持する。すると
\begin{equation}
|\psi\rangle
=
\alpha|\uparrow_z\rangle+\beta|\downarrow_z\rangle
\quad\longleftrightarrow\quad
\begin{pmatrix}
\alpha\\
\beta
\end{pmatrix},
\end{equation}
ここで
\begin{equation}
|\alpha|^2=P(\uparrow_z),
\qquad
|\beta|^2=P(\downarrow_z),
\qquad
|\alpha|^2+|\beta|^2=1.
\end{equation}
これは2次元の状態空間である。通常の空間が2次元であるという主張ではない。

\begin{classroomqa}
\lecturetimestamp{00:28:18}
\audiencequestion{物理空間が2次元であると言われたのですか。}
\lecturerresponse{いいえ。スピン状態空間が2次元である。位置と運動量は一時的に抑えられており、あたかも電子が壁に釘付けにされ、そのスピンだけが測定できるかのようである。}
\end{classroomqa}

この状態空間上の演算子は\(2\times2\)行列で表される。\(S_z\)基底を選んで\(S_z\)が対角になるようにすると便利であるが、対角性は物理的要請ではなく基底の選択である。

\begin{classroomqa}
\lecturetimestamp{00:29:46}
\audiencequestion{\(S_z\)は対角でなければならないのですか。}
\lecturerresponse{いいえ。便利さのために対角に選ばれている。他の行列表現も物理的に等価である。}
\end{classroomqa}

\begin{classroomqa}
\lecturetimestamp{00:30:08}
\audiencequestion{\(S_z\)は行列ですか。}
\lecturerresponse{\(S_z\)は可観測量であり、したがって演算子である。選ばれた有限次元基底では、その演算子は状態ベクトルに作用する行列で表される。}
\end{classroomqa}

特に有用な表現は
\begin{equation}
S_i=\frac12\sigma_i,
\end{equation}
ここで
\begin{equation}
\sigma_x=
\begin{pmatrix}
0&1\\
1&0
\end{pmatrix},
\quad
\sigma_y=
\begin{pmatrix}
0&-i\\
i&0
\end{pmatrix},
\quad
\sigma_z=
\begin{pmatrix}
1&0\\
0&-1
\end{pmatrix}.
\end{equation}
1/2の因子のない行列がパウリ行列であり、スピン演算子自体はその因子を含む。直接の掛け算で交換関係が確かめられる。表現は一意ではない：軸を回転させたり基底を変えたりすると、別の等価な三つ組が得られる。

これはスピン1/2角運動量の2成分表現である。

\subsection{座標軸に沿った固有状態}

固有値方程式を成分ごとに解くと有用である。\(\sigma_i\)の固有値は\(\pm1\)なので、\(S_i=\sigma_i/2\)の固有値は\(\pm\tfrac12\)である。

\(\sigma_z\)に対して、2つの方程式は
\begin{align}
\begin{pmatrix}
1&0\\
0&-1
\end{pmatrix}
\begin{pmatrix}
\alpha\\
\beta
\end{pmatrix}
&=
\begin{pmatrix}
\alpha\\
\beta
\end{pmatrix},
&
\begin{pmatrix}
1&0\\
0&-1
\end{pmatrix}
\begin{pmatrix}
\alpha\\
\beta
\end{pmatrix}
&=
-
\begin{pmatrix}
\alpha\\
\beta
\end{pmatrix}.
\end{align}

\begin{figure}[t]
\centering
\includegraphics[width=0.74\textwidth]{lecture_08_figure_02.png}
\caption{2成分スピノルに作用する\(\sigma_z\)の\(+1\)固有値方程式。}
\lectureframe{00:34:57}
\end{figure}

固有値\(+1\)に対して、行ごとの掛け算により
\begin{equation}
\begin{pmatrix}
1\cdot\alpha+0\cdot\beta\\
0\cdot\alpha-\beta
\end{pmatrix}
=
\begin{pmatrix}
\alpha\\
-\beta
\end{pmatrix}
=
\begin{pmatrix}
\alpha\\
\beta
\end{pmatrix}.
\end{equation}
上側の方程式は自動的に満たされ、下側の方程式から\(-\beta=\beta\)、したがって\(\beta=0\)が得られる。規格化は\(|\alpha|=1\)を要求する。無関係な全体位相を\(\alpha=1\)となるように選ぶと、
\begin{equation}
|+z\rangle=
\begin{pmatrix}
1\\
0
\end{pmatrix}.
\end{equation}

固有値\(-1\)に対して、
\begin{equation}
\begin{pmatrix}
\alpha\\
-\beta
\end{pmatrix}
=
\begin{pmatrix}
-\alpha\\
-\beta
\end{pmatrix}.
\end{equation}
今度は下側の方程式が自動的に満たされ、上側の方程式から\(\alpha=-\alpha\)、したがって\(\alpha=0\)が得られる。\(\beta=1\)と選ぶと
\begin{equation}
|-z\rangle=
\begin{pmatrix}
0\\
1
\end{pmatrix}.
\end{equation}

次に実験を変えて\(S_x\)を測定する。
\begin{equation}
\sigma_x
\begin{pmatrix}
\alpha\\
\beta
\end{pmatrix}
=
\begin{pmatrix}
0\cdot\alpha+\beta\\
\alpha+0\cdot\beta
\end{pmatrix}
=
\begin{pmatrix}
\beta\\
\alpha
\end{pmatrix},
\end{equation}
なので、\(+1\)の方程式は
\begin{equation}
\begin{pmatrix}
\beta\\
\alpha
\end{pmatrix}
=
\begin{pmatrix}
\alpha\\
\beta
\end{pmatrix}.
\end{equation}
となる。2つの成分方程式\(\beta=\alpha\)と\(\alpha=\beta\)は同じ条件である。規格化により
\begin{equation}
2|\alpha|^2=1,
\end{equation}
が得られるので、便利な選択は
\begin{equation}
|+x\rangle
=
\frac{1}{\sqrt2}
\begin{pmatrix}
1\\
1
\end{pmatrix}.
\end{equation}
である。

\begin{figure}[t]
\centering
\includegraphics[width=0.74\textwidth]{lecture_08_figure_03.png}
\caption{固有値\(+1\)に対して、\(\sigma_x\)は2つの成分を交換し、\(\alpha=\beta\)を要求する。}
\lectureframe{00:39:42}
\end{figure}

\begin{classroomqa}
\lecturetimestamp{00:40:13}
\audiencequestion{1/2の因子を含めると、固有値は1/2になりますか。}
\lecturerresponse{はい。パウリ行列\(\sigma_x\)は固有値\(\pm1\)をもち、一方\(S_x=\sigma_x/2\)は固有値\(\pm\tfrac12\)をもつ。}
\end{classroomqa}

固有値\(-1\)に対して、
\begin{equation}
\begin{pmatrix}
\beta\\
\alpha
\end{pmatrix}
=
-
\begin{pmatrix}
\alpha\\
\beta
\end{pmatrix},
\end{equation}
なので
\begin{equation}
\beta=-\alpha,
\qquad
\alpha=-\beta.
\end{equation}
これらもまた同じ方程式である。規格化の後、
\begin{equation}
|-x\rangle
=
\frac{1}{\sqrt2}
\begin{pmatrix}
1\\
-1
\end{pmatrix}.
\end{equation}
したがって、\(z\)アップと\(z\)ダウンの等しい重ね合わせは正の\(x\)軸に沿って偏極しており、相対的な負符号をもつ同じ重ね合わせは負の\(x\)軸に沿って偏極している。

\(y\)方向については、
\begin{equation}
\sigma_y
\begin{pmatrix}
\alpha\\
\beta
\end{pmatrix}
=
\begin{pmatrix}
-i\beta\\
i\alpha
\end{pmatrix}.
\end{equation}
規格化された2つの固有ベクトルは、全体位相を除いて
\begin{equation}
|+y\rangle
=
\frac{1}{\sqrt2}
\begin{pmatrix}
1\\
i
\end{pmatrix},
\qquad
|-y\rangle
=
\frac{1}{\sqrt2}
\begin{pmatrix}
1\\
-i
\end{pmatrix},
\end{equation}
である。どちらの状態も、\(z\)上向きに対して確率\(1/2\)を、\(z\)下向きに対して確率\(1/2\)を与える。これら2つの確率には現れない相対位相\(i\)が、\(y\)偏極と\(x\)偏極を区別する。線形結合は任意の方向に偏極したスピンも記述するが、その計算は後回しにする。

\begin{classroomqa}
\lecturetimestamp{00:44:22}
\audiencequestion{\(i\)成分の符号は誤った\(\sigma_y\)固有値に割り当てられてはいないか。}
\lecturerresponse{確かに符号が誤った固有値に割り当てられた可能性がある。ここではその割り当ては不確かであり、記憶から確定しているわけではない。}
\end{classroomqa}

直接の確認は
\begin{align}
\sigma_y\frac{1}{\sqrt2}
\begin{pmatrix}
1\\
i
\end{pmatrix}
&=
\frac{1}{\sqrt2}
\begin{pmatrix}
-i\,i\\
i
\end{pmatrix}
=
\frac{1}{\sqrt2}
\begin{pmatrix}
1\\
i
\end{pmatrix},\\
\sigma_y\frac{1}{\sqrt2}
\begin{pmatrix}
1\\
-i
\end{pmatrix}
&=
\frac{1}{\sqrt2}
\begin{pmatrix}
(-i)(-i)\\
i
\end{pmatrix}
=
-\frac{1}{\sqrt2}
\begin{pmatrix}
1\\
-i
\end{pmatrix}.
\end{align}
である。したがって\(+i\)は固有値\(+1\)に属し、\(-i\)は固有値\(-1\)に属する。\footnote{口頭での割り当ては何度も反転した後も未確定のままであり、ここに示した符号は表示された行列の積によって確定される。} スピン1/2のすべての状態は、任意の方向に偏極したすべての状態を含め、\(|+z\rangle\)と\(|-z\rangle\)の線形結合である。

\section{スピン1と位置依存スピノル}

スピン\(1\)に対して、可能な\(S_z\)の値は\(+1,0,-1\)である。したがって状態空間は3次元である。スピン\(1\)は整数スピンであり、ボソン統計に対応する。同じ角運動量交換関係を満たす3つの\(3\times3\)行列を求める。

3つの座標方向に関して対称な表現を以下に与える。\footnote{\(S_y\)の符号は\(00{:}48{:}50\)付近に見える行列に従い、また\([S_x,S_y]=iS_z\)によっても確定される。その符号を逆にすれば、要求される交換関係も逆になる。}
\begin{equation}
\begin{aligned}
S_x
&=
i
\begin{pmatrix}
0&0&0\\
0&0&1\\
0&-1&0
\end{pmatrix},
&
S_y
&=
i
\begin{pmatrix}
0&0&1\\
0&0&0\\
-1&0&0
\end{pmatrix},\\
S_z
&=
i
\begin{pmatrix}
0&1&0\\
-1&0&0\\
0&0&0
\end{pmatrix}.
\end{aligned}
\end{equation}
各行列は\(i\)倍の実反対称行列である。各行列は1つの\(1\)、1つの\(-1\)、およびそれ自身の座標に対応する消える行と列を含む。すなわち、1つ目は\(S_x\)、2つ目は\(S_y\)、3つ目は\(S_z\)に対応する。これは有用な記憶法を与える。この表現では\(S_z\)は対角ではない。

次のように置く。
\begin{equation}
v=
\begin{pmatrix}
\alpha\\
\beta\\
\gamma
\end{pmatrix}.
\end{equation}
すると
\begin{equation}
S_zv
=
\begin{pmatrix}
i\beta\\
-i\alpha\\
0
\end{pmatrix}.
\end{equation}
固有値\(+1\)に対して、
\begin{equation}
i\beta=\alpha,
\qquad
-i\alpha=\beta,
\qquad
0=\gamma.
\end{equation}
したがって\(\beta=-i\alpha\)および\(\gamma=0\)である。規格化により\(2|\alpha|^2=1\)となるので、
\begin{equation}
|m=+1\rangle
=
\frac{1}{\sqrt2}
\begin{pmatrix}
1\\
-i\\
0
\end{pmatrix}.
\end{equation}
固有値\(-1\)に対して、
\begin{equation}
i\beta=-\alpha,
\qquad
-i\alpha=-\beta,
\qquad
0=-\gamma,
\end{equation}
ゆえに\(\beta=i\alpha\)、\(\gamma=0\)、そして
\begin{equation}
|m=-1\rangle
=
\frac{1}{\sqrt2}
\begin{pmatrix}
1\\
i\\
0
\end{pmatrix}.
\end{equation}
最後に、固有値0の方程式から\(\alpha=\beta=0\)が得られ、\(\gamma\)は任意のままである。その規格化された解は
\begin{equation}
|m=0\rangle
=
\begin{pmatrix}
0\\
0\\
1
\end{pmatrix}.
\end{equation}
これ以外の固有値はない。これら3つの固有ベクトルが3次元空間を尽くす。因子\(1/\sqrt2\)は見かけ上のものではない。同じ係数は後に、この代数をアイソスピンに適用する際の定量的な実験予測にも現れる。

\begin{classroomqa}
\lecturetimestamp{00:54:10}
\audiencequestion{これらスピン1の固有ベクトルの複素成分はどのように読むべきか。}
\lecturerresponse{\(m=+1\)では成分は\(1/\sqrt2,-i/\sqrt2,0\)であり、\(m=-1\)では\(1/\sqrt2,i/\sqrt2,0\)であり、\(m=0\)では\(0,0,1\)である。}
\end{classroomqa}

他の行列の三つ組も同じ代数を満たし、物理的に等価である。上記の対称表現は三つの方向が同等に扱われるため有用であり、再び用いられる。

\subsection{位置の復元}

ここで粒子の運動を復元する。スピンは位置とも運動量とも可換であるから、どちらもスピン成分と同時に指定できる。位置と運動量は可換ではなく、異なるスピン成分も互いに可換ではない。

運動するスピン1/2粒子は位置依存スピノル
\begin{equation}
\psi(x)=
\begin{pmatrix}
\psi_\uparrow(x)\\
\psi_\downarrow(x)
\end{pmatrix}.
\end{equation}
で記述される。量
\begin{equation}
\rho_{\uparrow_z}(x)=|\psi_\uparrow(x)|^2,
\qquad
\rho_{\downarrow_z}(x)=|\psi_\downarrow(x)|^2
\end{equation}
は位置とスピンの結合確率密度である。その和
\begin{equation}
\rho(x)=\psi^\dagger(x)\psi(x)
=|\psi_\uparrow(x)|^2+|\psi_\downarrow(x)|^2,
\end{equation}
は位置密度である。粒子が\(x\)で検出されたことが分かっているなら、条件付き確率は
\begin{equation}
P(\uparrow_z\mid x)
=
\frac{|\psi_\uparrow(x)|^2}{\psi^\dagger(x)\psi(x)},
\qquad
P(\downarrow_z\mid x)
=
\frac{|\psi_\downarrow(x)|^2}{\psi^\dagger(x)\psi(x)},
\end{equation}
となる。ただし\(\rho(x)\neq0\)とする。したがって電子は2成分の空間波動関数で表され、局所的なスピン振幅は場所ごとに変化し得る。

運動するスピン1粒子に対して、
\begin{equation}
\phi(x)=
\begin{pmatrix}
\phi_1(x)\\
\phi_2(x)\\
\phi_3(x)
\end{pmatrix}.
\end{equation}
と書く。空間の各点で、3つの成分は上で用いたのと同じスピン空間のベクトルをなす。スピン1粒子は整数スピンをもち、ボソンである。

\begin{classroomqa}
\lecturetimestamp{01:00:10}
\audiencequestion{\(\phi_1(x)\)自体がスピン1の結果\(m=+1\)に対する確率なのか。}
\lecturerresponse{いいえ。それは選ばれた三成分基底における振幅である。確定したスピン結果を得るには、対応する規格化された固有ベクトルのブラを作り、場をそれに射影し、その絶対値の二乗を取る。}
\end{classroomqa}

\begin{figure}[t]
\centering
\includegraphics[width=0.82\textwidth]{lecture_08_figure_04.png}
\caption{規格化された\(\sigma_y\)固有ベクトルと、3つのスピン1/2行列。}
\lectureframe{01:00:12}
\end{figure}

例えば、
\begin{equation}
\begin{aligned}
|m=+1\rangle
&=
\frac{1}{\sqrt2}
\begin{pmatrix}
1\\
-i\\
0
\end{pmatrix}
\\
\Longrightarrow\quad\langle m=+1|
&=
\frac{1}{\sqrt2}
\begin{pmatrix}
1&i&0
\end{pmatrix}.
\end{aligned}
\end{equation}
したがって射影振幅は
\begin{align}
\langle m=+1|\phi(x)\rangle
&=
\frac{1}{\sqrt2}
\begin{pmatrix}
1&i&0
\end{pmatrix}
\begin{pmatrix}
\phi_1(x)\\
\phi_2(x)\\
\phi_3(x)
\end{pmatrix}\\
&=
\frac{\phi_1(x)+i\phi_2(x)}{\sqrt2}.
\end{align}
同じ規則を三つの固有ベクトルすべてに適用すると、スピン分解された密度が得られる。\footnote{これらの式は、\(\phi_1\) と \(\phi_2\) が実数である必要はないという明示的な補正を経た後の射影代数を完成させる。それらを実数とみなすと、\(m=+1\) 密度と \(m=-1\) 密度を区別する干渉が誤って消去されてしまう。}
\begin{align}
\rho_{+1}(x)
&=
\frac12\left|\phi_1(x)+i\phi_2(x)\right|^2,\\
\rho_{-1}(x)
&=
\frac12\left|\phi_1(x)-i\phi_2(x)\right|^2,\\
\rho_0(x)
&=
|\phi_3(x)|^2.
\end{align}
これらは次を満たす。
\begin{equation}
\begin{aligned}
\rho_{+1}(x)+\rho_{-1}(x)+\rho_0(x)
&=
|\phi_1(x)|^2+|\phi_2(x)|^2
\\
&\quad{}+|\phi_3(x)|^2
\\
&=
\phi^\dagger(x)\phi(x).
\end{aligned}
\end{equation}
粒子が \(x\) で検出されたという条件が課されるならば、
\begin{equation}
P(m\mid x)
=
\frac{\rho_m(x)}{\phi^\dagger(x)\phi(x)},
\end{equation}
ただし \(\phi^\dagger(x)\phi(x)\neq0\) とする。一般的な測定規則は常に同じである。すなわち、望む規格化された固有ベクトルのブラを作り、実際の状態をそれへ射影し、その絶対値の二乗を取る。得られる密度はその位置依存性を保つ。

より高いスピンの行列はより複雑であるが、原理は変わらない。

\begin{classroomqa}
\lecturetimestamp{01:07:31}
\audiencequestion{スピン3/2の状態空間はどれくらいの大きさか。}
\lecturerresponse{それは4つの状態をもち、\(m=\tfrac32,\tfrac12,-\tfrac12,-\tfrac32\)である。スピン2は5つの状態をもつ。}
\end{classroomqa}

\begin{classroomqa}
\lecturetimestamp{01:07:55}
\audiencequestion{なぜ議論には\(x,y,z\)が現れるのに時間方向は現れないのか。この扱いは非相対論的か。}
\lecturerresponse{今のところは非相対論的である。スピンは質量と同様、粒子の静止系で最も簡単に特徴づけられ、別の慣性系での記述は適切なローレンツ変換によって得られる。その相対論的な段階は後回しにする。}
\end{classroomqa}

ディラック方程式は相対論的な枠組みを与え、その4成分波動関数がなぜスピン二重項を含むのかを説明する。

\section{ディラック行列と4成分スピノル}

4成分ディラック波動関数は次を満たす。
\begin{equation}
i\frac{\partial\Psi}{\partial t}
=
-i\alpha_i\frac{\partial\Psi}{\partial x_i}
+m\beta\Psi,
\qquad i=1,2,3,
\end{equation}
ここで繰り返される空間添字は和をとる。したがって微分項は
\begin{equation}
\alpha_1\frac{\partial\Psi}{\partial x}
+\alpha_2\frac{\partial\Psi}{\partial y}
+\alpha_3\frac{\partial\Psi}{\partial z}.
\end{equation}
を意味する。平面波に対してこれは代数方程式
\begin{equation}
\omega\Psi=(\alpha_i k_i+m\beta)\Psi.
\end{equation}
となる。

\begin{classroomqa}
\lecturetimestamp{01:11:32}
\audiencequestion{\(m\)の2回の出現は同じものを表すのか。}
\lecturerresponse{いいえ。一方は空間添字として用いられ、他方は粒子の質量である。空間添字を\(i\)と改名する。\(i\)が2回現れるときは\(1\)から\(3\)まで和をとる。}
\end{classroomqa}

振動数演算子を二乗すると、相対論的関係
\begin{equation}
\omega^2=\mathbf{k}^2+m^2.
\end{equation}
が再現されなければならない。したがって二乗の項は、各 \(i\) について次を要求する。
\begin{equation}
\alpha_i^2=I_4,
\qquad
\beta^2=I_4.
\end{equation}

\begin{classroomqa}
\lecturetimestamp{01:13:01}
\audiencequestion{\(\beta\)は\(\Psi\)の4つの成分のそれぞれで異なるのか。}
\lecturerresponse{いいえ。4成分スピノル全体に作用する1つの行列\(\beta\)があり、それは\(\beta^2=I_4\)を満たす。}
\end{classroomqa}

混合項は消えなければならない。異なる空間方向の間で、またすべての\(\alpha_i\)と\(\beta\)の間で、これは
\begin{equation}
\alpha_i\alpha_j+\alpha_j\alpha_i=0
\quad (i\neq j),
\qquad
\alpha_i\beta+\beta\alpha_i=0.
\end{equation}
を与える。等価に、
\begin{equation}
\{\alpha_i,\alpha_j\}=2\delta_{ij}I_4,
\quad
\{\alpha_i,\beta\}=0,
\quad
\{\beta,\beta\}=2I_4.
\end{equation}
各行列は異なるすべての行列と反交換し、自身との反交換子は単位行列の2倍である。

\(\beta\) がなければ、三つのパウリ行列がすでに、二乗すると単位行列になり、互いに反交換する三つの \(2\times2\) 行列を与えるであろう。質量項はそのような第四の行列を要求する。要求される性質をもつ第四の \(2\times2\) 行列は存在せず、\(3\times3\) 行列もこの代数を解かない。したがって最小の表現は \(4\times4\) である。

\(4\times4\)行列は4つの\(2\times2\)ブロックに分けることができる。1つの便利な表現では、
\begin{equation}
\alpha_i=
\begin{pmatrix}
\sigma_i&0\\
0&-\sigma_i
\end{pmatrix},
\qquad i=1,2,3.
\end{equation}

\begin{figure}[t]
\centering
\includegraphics[width=0.68\textwidth]{lecture_08_figure_05.png}
\caption{\(2\times2\)ブロックで書かれたディラック行列\(\alpha_i\)であり、対角線上に\(\sigma_i\)と\(-\sigma_i\)をもつ。}
\lectureframe{01:17:23}
\end{figure}

\begin{classroomqa}
\lecturetimestamp{01:17:41}
\audiencequestion{先の表現では\(\sigma_i\)ブロックは逆の配置だったのか。}
\lecturerresponse{はい。別の表現ではそれらは非対角に置かれていた。2つの配置は等価であり、相似変換で関係づけられる。}
\end{classroomqa}

現在の表現では残りの行列は
\begin{equation}
\beta=
\begin{pmatrix}
0&I_2\\
I_2&0
\end{pmatrix},
\qquad
I_2=
\begin{pmatrix}
1&0\\
0&1
\end{pmatrix}.
\end{equation}
である。ブロック積により、\(4\times4\)代数の確認は\(2\times2\)行列の通常の積に帰着する。

\begin{classroomqa}
\lecturetimestamp{01:19:17}
\audiencequestion{先の表現では\(\beta\)に異なる符号または対角形が用いられていたのか。}
\lecturerresponse{それは異なるが物理的に等価な基底を用いていた。基底の変更は\(\Psi\)の4つの成分を混合し、物理を変えずに行列の見かけを変える。}
\end{classroomqa}

\(\Psi\) の4つの成分は、時空の4つの座標ではない。これらは、この反交換代数が作用できる最小のベクトル空間における座標である。2つの異なる二重構造が存在する。\(2\times2\) ブロックへの分割と、各ブロック内の2成分パウリ構造である。以下で確立される解釈を先取りすると、4つの可能性は次のように整理できる。
\begin{equation}
\begin{array}{c|cc}
& S_z=+\tfrac12 & S_z=-\tfrac12\\ \hline
\omega>0 & (\omega>0,+\tfrac12) & (\omega>0,-\tfrac12)\\
\omega<0 & (\omega<0,+\tfrac12) & (\omega<0,-\tfrac12)
\end{array}.
\end{equation}
下の行は負振動数の電子状態からなる。それは陽電子の行ではない。陽電子は、それらの状態の満たされた集合における空孔として解釈されるであろう。

\section{静止系におけるエネルギーの倍化と陽電子}

二重のエネルギー構造は、ゼロ運動量で最も容易に露呈する。\(\mathbf{k}=0\) のとき、平面波状態は空間的に一定であり、空間微分の項が消える。静止系の方程式\footnote{\(01{:}23{:}38\) における枠は、質量項に一時的な余分な因子 \(i\) を保持している。\(01{:}24{:}05\) における口頭の訂正がそれを取り除く。} は
\begin{equation}
i\frac{\partial\Psi}{\partial t}=m\beta\Psi.
\end{equation}

\begin{figure}[t]
\centering
\includegraphics[width=0.76\textwidth]{lecture_08_figure_06.png}
\caption{ゼロ運動量では空間微分が消え、訂正された静止系の方程式 \(i\partial_t\Psi=m\beta\Psi\) が残る。}
\lectureframe{01:23:38}
\end{figure}

4成分の列を、上と下の2成分の対象に分割する。
\begin{equation}
\Psi=
\begin{pmatrix}
\psi_+\\
\psi_-
\end{pmatrix}.
\end{equation}

\begin{classroomqa}
\lecturetimestamp{01:24:47}
\audiencequestion{ラベル \(\psi_+\) と \(\psi_-\) は、すでに正エネルギーと負エネルギーを意味しているのか。}
\lecturerresponse{いいえ。この段階では、それらは上と下の2成分ブロックの単なる名前である。それらのエネルギーの組み合わせはまだ見出されていない。}
\end{classroomqa}

\begin{classroomqa}
\lecturetimestamp{01:25:10}
\audiencequestion{分割された4成分の列は何を表すのか。}
\lecturerresponse{それは同じ4成分ディラックスピノルであり、行列のブロック形式に合うように整理されている。それぞれの半分は2成分の対象であり、そのような2つの対象が元の4成分の列を作る。}
\end{classroomqa}

行列 \(\beta\) は2つのブロックを交換する。
\begin{equation}
\beta
\begin{pmatrix}
\psi_+\\
\psi_-
\end{pmatrix}
=
\begin{pmatrix}
\psi_-\\
\psi_+
\end{pmatrix}.
\end{equation}
したがって静止系の方程式は次のようになる。
\begin{equation}
i\frac{\partial\psi_+}{\partial t}=m\psi_-,
\qquad
i\frac{\partial\psi_-}{\partial t}=m\psi_+.
\end{equation}

2つの自然な提案がある。もう一度時間微分を取れば系は分離するが、2つの方程式を加えたり引いたりすれば直ちに対角化される。実際、
\begin{align}
i\frac{\partial}{\partial t}(\psi_++\psi_-)
&=m(\psi_++\psi_-),\\
i\frac{\partial}{\partial t}(\psi_+-\psi_-)
&=-m(\psi_+-\psi_-).
\end{align}
次のように定義すると
\begin{equation}
\chi_+=\psi_++\psi_-,
\qquad
\chi_-=\psi_+-\psi_-,
\end{equation}
2つの独立な方程式が得られる。
\begin{equation}
i\partial_t\chi_+=m\chi_+,
\qquad
i\partial_t\chi_-=-m\chi_-.
\end{equation}
確定した振動数依存性 \(e^{-i\omega t}\) に対して、和は
\begin{equation}
\omega=+m,
\end{equation}
を持ち、一方、差は
\begin{equation}
\omega=-m.
\end{equation}
を持つ。これはまさに、以前の1次元の例で見出された正エネルギーと負エネルギーの倍化である。すなわち、相対論的方程式は再び正エネルギーと負エネルギーの電子解を持つ。したがって、ブロックの倍化は正エネルギー枝と負エネルギー枝に分かれる。純粋な静止系状態は次のように書ける。
\begin{equation}
\Psi_{\omega=+m}=
\begin{pmatrix}
\xi\\
\xi
\end{pmatrix},
\qquad
\Psi_{\omega=-m}=
\begin{pmatrix}
\xi\\
-\xi
\end{pmatrix},
\end{equation}
ここで \(\xi\) は依然として任意の2成分の対象である。

和と差は一方の二重構造を片付けたが、\(\xi\) の2つの成分が残る。パウリ行列はこの残りの対に作用する。これはスピンに関連する対であり、以下で議論する回転と保存則の検査に従う。

\begin{classroomqa}
\lecturetimestamp{01:32:23}
\audiencequestion{スピンを外側の二重構造として、エネルギーの符号を内側の二重構造として扱えるか。}
\lecturerresponse{はい。2つの因子の順序は等価である。どちらの添字を先に書くかを入れ替えてもよい。その変更は4つの成分の並べ替えにすぎない。}
\end{classroomqa}

負振動数枝はディラックの海の解釈へと導く。負振動数には負の電子エネルギーが割り当てられる。すべての負エネルギー電子状態を満たすことは、排他原理が許す限りエネルギーを下げる。なぜなら、各完全なフェルミオン状態を占めることができる電子は1つだけだからである。完全に満たされた負エネルギーの海は、そのとき真空と呼ばれる。

この満たされた海から1つの電子を正エネルギー状態へ昇格させる。その結果は、正エネルギーの電子と、占められていない負エネルギー状態を含む。負エネルギーを取り除くことはエネルギーを上げるので、空孔は正の励起エネルギーを持つ。電子が欠けているため、それは反対の電荷も運ぶ。その空孔は陽電子である。

\begin{classroomqa}
\lecturetimestamp{01:32:57}
\audiencequestion{陽電子の解釈をもう一度説明できるか。}
\lecturerresponse{2つのブロックの和は正振動数を持ち、差は負振動数を持つ。負振動数の解は負エネルギーの電子状態であり、真空はそれらすべてを満たすことによって得られる。満たされた1つの電子を正エネルギーへ昇格させると、正エネルギーの電子と空孔の両方が残る。その正エネルギーで反対の電荷を持つ空孔が陽電子である。}
\end{classroomqa}

したがって、この代数は同時に二つの帰結を生み出した。すなわち、粒子--反粒子構造となる正エネルギー枝と負エネルギー枝、そしてスピン1/2に対応する残存する2成分構造である。注目すべき点は、その両方が、相対論的反交換関係を満たしうる最小の行列から現れたことである。この数学的構成がスピンをもつ粒子と反粒子枝を生み出したとき、スピンはごく最近になって経験的性質として現れたばかりであった。その歴史的状況において、この結果は劇的であった。

\begin{classroomqa}
\lecturetimestamp{01:35:14}
\audiencequestion{負エネルギーの電子それ自体が陽電子なのですか。}
\lecturerresponse{いいえ。負エネルギーの電子状態は真空において詰まっている。陽電子とは、その詰まった集合から電子が一つ欠けていること、すなわち空孔であり、負エネルギーの電子ではない。}
\end{classroomqa}

空孔は遷移に関与しうる。陽電子が正のエネルギーをもつのは、それが負エネルギーの電子の不在だからである。正エネルギーの電子が空孔を埋めれば、電子と陽電子は対消滅するが、そのエネルギーは運び去られねばならず、通常は光子によって運ばれる。海の言葉では、電子が空の負エネルギー状態を埋めたことになる。同じエネルギーの放出は、原子内の励起電子がより低い準位へ落ちる過程として描くこともできる。下向きの遷移は光子を放出する。

\begin{classroomqa}
\lecturetimestamp{01:36:39}
\audiencequestion{なぜ、より高い負エネルギーの電子は空孔へ落ちないのですか。}
\lecturerresponse{その問いは二つの場合を混同している。空孔を埋める正エネルギーの電子は電子--陽電子対消滅であり、光子がエネルギーを運び去る。負エネルギーの海の電子による遷移は別の過程であり、以下で扱う。}
\end{classroomqa}

運動量ゼロの計算それ自体は、エネルギーの符号を分離した後に二重縮退を確立した。二つの状態の縮退は、それ自体では、それらの状態が角運動量をもつことの証明ではない。

\begin{classroomqa}
\lecturetimestamp{01:37:44}
\audiencequestion{運動量ゼロの計算は、残った一対の状態がスピンであることをどのように証明するのですか。}
\lecturerresponse{その一対は上と下と呼ばれてよく、パウリ行列は完全な方程式に現れるが、それだけではまだ証明ではない。その一対が空間回転の下でどのように変換するかを示すか、あるいはそれにスピンを割り当てることで相互作用や衝突において全角運動量が保存されることを示さねばならない。その証明は後回しにする。}
\end{classroomqa}

正エネルギーの静止状態では、上側と下側のブロックは等しく、
\begin{equation}
\Psi_{\omega=+m}=
\begin{pmatrix}
\xi\\
\xi
\end{pmatrix}.
\end{equation}
\(\xi\) が角運動量の二重項であることを確立するには、いま述べた回転または保存の議論が必要であり、二つの成分を数えるだけでは従わない。

\begin{classroomqa}
\lecturetimestamp{01:40:46}
\audiencequestion{たとえそのエネルギーが異なっていても、別の電子が陽電子と対消滅できますか。}
\lecturerresponse{はい。電子は個々の同一性を保持しない。利用可能などちらの電子でも空孔を埋めることができ、初期のエネルギーと運動量の違いは、放出される光子の数とエネルギーに反映される。}
\end{classroomqa}

空孔描像は、負エネルギーの海の内部に完全に留まる遷移も説明する。占有された負エネルギー状態が空孔へ移ったとしよう。元の空孔は埋まるが、その電子が来た状態は今や空である。空孔が移動したのである。等価に、陽電子はそのエネルギーを変え、放射を放出したのである。

\begin{classroomqa}
\lecturetimestamp{01:42:12}
\audiencequestion{なぜ、隣り合う負エネルギーの準位が次々と空孔へ落ちることができないのですか。}
\lecturerresponse{それらは空孔を一つの準位から別の準位へ移すことができる。電子が古い空孔へ落ちると、その以前のエネルギーに新しい空孔が残るので、初期の陽電子は異なるエネルギーの陽電子になる。エネルギー差は光子によって運ばれる。}
\end{classroomqa}

\begin{classroomqa}
\lecturetimestamp{01:44:39}
\audiencequestion{なぜこれは際限のない下向きの放射カスケードを生じないのですか。}
\lecturerresponse{エネルギーと運動量が同時に保存されねばならないからである。それら二つの保存則は無制限の過程を許さない。同じ問題は通常の自由電子にも生じる。それは真空中で光子を放出して単にエネルギーを下げ続けることはできない。}
\end{classroomqa}
